% Flavor catalog per-process template.
% process_id: E009
% family: edm_neutrino
% owner_agent_id: PKA-E009
% writer_agent_id: null
% checker_agent_id: null
% opus_signoff_id: null
% Canonical metadata sidecar: E009.yaml
% Status is controlled by the YAML sidecar status_history, not this comment.

\section{E009: Weinberg three-gluon operator}

\subsection*{Process}
\textbf{Standard notation:} \(w\) or \(C_{\tilde G}\), multiplying the
CP-odd gluonic operator schematically
\(f^{abc}\tilde G^a_{\mu\nu}G^{b\,\nu\rho}G^c_{\rho}{}^{\mu}\).  This entry
covers EDM constraints on the flavor-diagonal Weinberg three-gluon operator,
not the lepton-number Weinberg operator for neutrino masses.

\subsection*{PDG / equivalent value(s)}
There is no directly measured PDG value for \(w\).  The measured anchor is the
PDG Live 2026 neutron EDM limit
\[
  |d_n| < 1.8\times 10^{-26}\ e\,\mathrm{cm}\quad(90\%~\mathrm{CL}),
\]
based on Abel et al. 2020, with
\[
  d_n=(0.0\pm1.1_{\rm stat}\pm0.2_{\rm sys})\times10^{-26}\ e\,\mathrm{cm}.
\]
The local snapshots are
\texttt{flavor\_catalog/references/E009/pdg2026\_neutron\_edm\_datablock.txt}
and \texttt{flavor\_catalog/references/E009/abel2020\_neutron\_edm\_arxiv2001\_11966.txt}.
Theoretical normalizations and one-source-at-a-time Wilson-coefficient bounds
are recorded only in the YAML auxiliary blocks.

\subsection*{Relevance to the RS / anarchic-flavor pipeline}
Anarchic warped or partial-compositeness models carry new CP phases in the
colored sector.  Integrating out KK fermions, heavy quarks, Higgs-sector states,
or colored resonances can generate quark chromo-EDMs and finite heavy-threshold
contributions to the Weinberg operator.  The operator is therefore a compact
EDM-adjacent diagnostic of colored CP violation that complements E008.

\subsection*{Post-2008 developments}
Relative to the arXiv:0804.1954 RS-flavor baseline, the experimental neutron
anchor moved to the Abel 2020/PDG 2026 limit above.  Composite/warped dipole
analyses after 2008 emphasized that heavy-quark CEDM running generates a
finite threshold correction to the three-gluon Weinberg operator
(\texttt{KoenigNeubertStraub:CompositeDipoles2014}).  The standard
Pospelov--Ritz paper-era estimate gives
\(|d_n(w)|\simeq e\,22~\mathrm{MeV}\,w(1~\mathrm{GeV})\), which implies the
central benchmark \(|w(1~\mathrm{GeV})|<4.1\times10^{-11}\ \mathrm{GeV}^{-2}\)
from the current neutron limit if all other CP-odd sources vanish.  A later
sum-rule update gives, in its own \(O_6\) convention,
\((d_n/e)_{O_6}=74(1\pm0.5)~\mathrm{MeV}\), corresponding to the central
benchmark \(|C_6|<1.2\times10^{-11}\ \mathrm{GeV}^{-2}\)
(\texttt{HaischHala:WeinbergSumRules2019}).  Lattice work using gradient flow
now explicitly studies the dimension-six gluonic CP-violating matrix element
and its operator mixing (\texttt{Bhattacharya:WeinbergLattice2022}).

\subsection*{Constraint validity and model dependence}
The neutron EDM limit is robust.  The Weinberg-operator bound is not: it
depends on operator normalization, RG running, heavy-quark thresholds,
hadronic matrix elements, and assumptions about cancellations with
\(\bar\theta\), quark EDMs, qCEDMs, semileptonic terms, and four-quark
operators.  In global EDM language, the same neutron and diamagnetic data
probe more CP-odd sources than there are independent low-energy observables.

\subsection*{Code coverage in this repo}
\textbf{NO.}  Required greps over \texttt{quarkConstraints/}, \texttt{qcd/},
\texttt{flavorConstraints/}, \texttt{neutrinos/}, \texttt{yukawa/},
\texttt{warpConfig/}, \texttt{solvers/}, \texttt{scanParams/}, and
\texttt{tests/} found no Weinberg, neutron-EDM, or CP-odd gluonic operator
constraint.  The only nearby dipole implementation is unrelated
\(\mu\to e\gamma\) code in \texttt{flavorConstraints/muToEGamma.py:3},
\texttt{:21}, and \texttt{:81}.

\subsection*{Implementation difficulty}
\textbf{HIGH.}  A live E009 constraint needs new CP-odd colored operators,
RS loop or threshold matching, QCD running and mixing into other EDM sources,
and a chosen hadronic matrix-element convention.  The existing \(\Delta F=2\)
basis and muon LFV dipole path do not cover this observable.

\subsection*{Key references}
Process-local keys before bibliography consolidation:
\texttt{PDG2026:NeutronEDM}, \texttt{Abel:NeutronEDM2020},
\texttt{Weinberg:ThreeGluon1989}, \texttt{PospelovRitz:EDMReview2005},
\texttt{ChuppRamseyMusolf:EDMGlobal2014},
\texttt{KoenigNeubertStraub:CompositeDipoles2014},
\texttt{HaischHala:WeinbergSumRules2019},
\texttt{Bhattacharya:WeinbergLattice2022}, and
\texttt{CsakiFalkowskiWeiler:RSFlavor2008}.
